\documentclass[12pt]{article}

% Idioma y codificación
\usepackage[spanish, es-tabla, es-notilde]{babel}       %es-tabla para que se titule "Tabla"
\usepackage[utf8]{inputenc}

% Márgenes
\usepackage[a4paper,top=3cm,bottom=2.5cm,left=3cm,right=3cm]{geometry}

% Comentarios de bloque
\usepackage{verbatim}

% Paquetes de links
\usepackage[hidelinks]{hyperref}    % Permite enlaces
\usepackage{url}                    % redirecciona a la web

% Más opciones para enumeraciones
\usepackage{enumitem}

% Personalizar la portada
\usepackage{titling}

% Paquetes de tablas
\usepackage{multirow}

% Para añadir el símbolo de euro
\usepackage{eurosym}

% Paquete para subrayar con diferentes estilos
\usepackage[normalem]{ulem}


%------------------------------------------------------------------------

%Paquetes de figuras
\usepackage{caption}
\usepackage{subcaption} % Figuras al lado de otras
\usepackage{float}      % Poner figuras en el sitio indicado H.


% Paquetes de imágenes
\usepackage{graphicx}       % Paquete para añadir imágenes
\usepackage{transparent}    % Para manejar la opacidad de las figuras

% Paquete para usar colores
\usepackage[dvipsnames, table, xcdraw]{xcolor}
\usepackage{pagecolor}      % Para cambiar el color de la página

% Habilita tamaños de fuente mayores
\usepackage{fix-cm}

% Para los gráficos
\usepackage{tikz}
\usepackage{forest}

% Para poder situar los nodos en los grafos
\usetikzlibrary{positioning}

% Para evitar problemas entre babel y tikz
\usetikzlibrary{babel}


%------------------------------------------------------------------------

% Paquetes de matemáticas
\usepackage{mathtools, amsfonts, amssymb, mathrsfs}
\usepackage[makeroom]{cancel}     % Simplificar tachando
\usepackage{polynom}    % Divisiones y Ruffini
\usepackage{units} % Para poner fracciones diagonales con \nicefrac

\usepackage{pgfplots}   %Representar funciones
\pgfplotsset{compat=1.18}  % Versión 1.18

\usepackage{tikz-cd}    % Para usar diagramas de composiciones
\usetikzlibrary{calc}   % Para usar cálculo de coordenadas en tikz

%Definición de teoremas, etc.
\usepackage{amsthm}
%\swapnumbers   % Intercambia la posición del texto y de la numeración

\theoremstyle{plain}

\makeatletter
\@ifclassloaded{article}{
  % --- Configuración para ARTICLE ---
  \newtheorem{teo}{Teorema}[section]
  \theoremstyle{definition}
  \newtheorem{definicion}{Definición}[section]
  \newtheorem{ejercicio}{Ejercicio}[section]

  \renewcommand{\theteo}{%
    \ifnum\value{section}=0 \arabic{teo}\else \thesection.\arabic{teo}\fi
  }
  \renewcommand{\thedefinicion}{%
    \ifnum\value{section}=0 \arabic{definicion}\else \thesection.\arabic{definicion}\fi
  }
  \renewcommand{\theejercicio}{%
    \ifnum\value{section}=0 \arabic{ejercicio}\else \thesection.\arabic{ejercicio}\fi
  }
}{
  % --- Configuración para BOOK / REPORT ---
  \newtheorem{teo}{Teorema}[chapter]
  \theoremstyle{definition}
  \newtheorem{definicion}{Definición}[chapter]
  \newtheorem{ejercicio}{Ejercicio}[section]

  % Si no hay sección: Capítulo.Teorema (ej. 2.1)
  % Si hay sección:    Capítulo.Sección.Teorema (ej. 2.3.1)
  \renewcommand{\theteo}{%
    \ifnum\value{section}=0
      \thechapter.\arabic{teo}%
    \else
      \thesection.\arabic{teo}% (\thesection ya contiene el número de capítulo)
    \fi
  }
  \renewcommand{\thedefinicion}{%
    \ifnum\value{section}=0
      \thechapter.\arabic{definicion}%
    \else
      \thesection.\arabic{definicion}%
    \fi
  }
  \renewcommand{\theejercicio}{%
    \ifnum\value{section}=0
      \arabic{ejercicio}%
    \else
      \thesection.\arabic{ejercicio}% Formato por defecto: ej. 2.2.1
    \fi
  }
}
\makeatother

% Teoremas dependientes o con contador compartido
\newtheorem{coro}{Corolario}[teo]
\newtheorem{prop}[teo]{Proposición}
\newtheorem{lema}[teo]{Lema}
\newtheorem*{lema*}{Lema}

\theoremstyle{remark}
\newtheorem*{observacion}{Observación}

\theoremstyle{definition}
\newtheorem*{notacion}{Notación}
\newtheorem*{ejemplo}{Ejemplo}
\newtheorem*{ejercicio*}{Ejercicio}


% \renewcommand\qedsymbol{$\blacksquare$}         % Cambiar símbolo QED
%------------------------------------------------------------------------

% Paquetes para encabezados
\usepackage{fancyhdr}
\pagestyle{fancy}
\fancyhf{}

\newcommand{\helv}{ % Modificación tamaño de letra
\fontfamily{}\fontsize{12}{12}\selectfont}
\setlength{\headheight}{15pt} % Amplía el tamaño del índice


%\usepackage{lastpage}   % Referenciar última pag   \pageref{LastPage}
%\fancyfoot[C]{%
%  \begin{minipage}{\textwidth}
%    \centering
%    ~\\
%    \thepage\\
%    \href{https://losdeldgiim.github.io/}{\texttt{\footnotesize losdeldgiim.github.io}}
%  \end{minipage}
%}
\fancyfoot[C]{\thepage}
\fancyfoot[R]{\href{https://losdeldgiim.github.io/}{\texttt{\footnotesize losdeldgiim.github.io}}}

%------------------------------------------------------------------------

% Conseguir que no ponga "Capítulo 1". Sino solo "1."
\makeatletter
\@ifclassloaded{book}{
  \renewcommand{\chaptermark}[1]{\markboth{\thechapter.\ #1}{}} % En el encabezado
    
  \renewcommand{\@makechapterhead}[1]{%
  \vspace*{50\p@}%
  {\parindent \z@ \raggedright \normalfont
    \ifnum \c@secnumdepth >\m@ne
      \huge\bfseries \thechapter.\hspace{1em}\ignorespaces
    \fi
    \interlinepenalty\@M
    \Huge \bfseries #1\par\nobreak
    \vskip 40\p@
  }}
}
\makeatother

%------------------------------------------------------------------------
% Paquetes de cógido
\usepackage{minted}
\renewcommand\listingscaption{Código fuente}

\usepackage{fancyvrb}
% Personaliza el tamaño de los números de línea
\renewcommand{\theFancyVerbLine}{\small\arabic{FancyVerbLine}}

% Estilo para C++
\newminted{cpp}{
    frame=lines,
    framesep=2mm,
    baselinestretch=1.2,
    linenos,
    escapeinside=||
}

% para minted
\definecolor{LightGray}{rgb}{0.95,0.95,0.92}
\setminted{
    linenos=true,
    stepnumber=5,
    numberfirstline=true,
    autogobble,
    breaklines=true,
    breakautoindent=true,
    breaksymbolleft=,
    breaksymbolright=,
    breaksymbolindentleft=0pt,
    breaksymbolindentright=0pt,
    breaksymbolsepleft=0pt,
    breaksymbolsepright=0pt,
    fontsize=\footnotesize,
    bgcolor=LightGray,
    numbersep=10pt
}


\usepackage{listings} % Para incluir código desde un archivo

\renewcommand\lstlistingname{Código Fuente}
\renewcommand\lstlistlistingname{Índice de Códigos Fuente}

% Definir colores
\definecolor{vscodepurple}{rgb}{0.5,0,0.5}
\definecolor{vscodeblue}{rgb}{0,0,0.8}
\definecolor{vscodegreen}{rgb}{0,0.5,0}
\definecolor{vscodegray}{rgb}{0.5,0.5,0.5}
\definecolor{vscodebackground}{rgb}{0.97,0.97,0.97}
\definecolor{vscodelightgray}{rgb}{0.9,0.9,0.9}

% Configuración para el estilo de C similar a VSCode
\lstdefinestyle{vscode_C}{
  backgroundcolor=\color{vscodebackground},
  commentstyle=\color{vscodegreen},
  keywordstyle=\color{vscodeblue},
  numberstyle=\tiny\color{vscodegray},
  stringstyle=\color{vscodepurple},
  basicstyle=\scriptsize\ttfamily,
  breakatwhitespace=false,
  breaklines=true,
  captionpos=b,
  keepspaces=true,
  numbers=left,
  numbersep=5pt,
  showspaces=false,
  showstringspaces=false,
  showtabs=false,
  tabsize=2,
  frame=tb,
  framerule=0pt,
  aboveskip=10pt,
  belowskip=10pt,
  xleftmargin=10pt,
  xrightmargin=10pt,
  framexleftmargin=10pt,
  framexrightmargin=10pt,
  framesep=0pt,
  rulecolor=\color{vscodelightgray},
  backgroundcolor=\color{vscodebackground},
}

%------------------------------------------------------------------------
% La etiqueda IA sirve para indicar a partir de donde le pasamos el preambulo a la IA para el escaneo
% IA

% Comandos definidos
\newcommand{\bb}[1]{\mathbb{#1}}
\newcommand{\cc}[1]{\mathcal{#1}}

% I prefer the slanted \leq
\let\oldleq\leq % save them in case they're every wanted
\let\oldgeq\geq
\renewcommand{\leq}{\leqslant}
\renewcommand{\geq}{\geqslant}

% Si y solo si
\newcommand{\sii}{\iff}

% MCD y MCM
\DeclareMathOperator{\mcd}{mcd}
\DeclareMathOperator{\mcm}{mcm}

% Signo
\DeclareMathOperator{\sgn}{sgn}

% Imagen
\DeclareMathOperator{\Ima}{Im}

% Aplicación identidad
\newcommand{\id}{\mathrm{Id}}

% Letras griegas
\newcommand{\eps}{\epsilon}
\newcommand{\veps}{\varepsilon}
\newcommand{\lm}{\lambda}

\newcommand{\ol}{\overline}
\newcommand{\ul}{\underline}
\newcommand{\wt}{\widetilde}
\newcommand{\wh}{\widehat}

\let\oldvec\vec
\renewcommand{\vec}{\overrightarrow}

% Derivadas parciales
\newcommand{\del}[2]{\frac{\partial #1}{\partial #2}}
\newcommand{\Del}[3]{\frac{\partial^{#1} #2}{\partial #3^{#1}}}
\newcommand{\deld}[2]{\dfrac{\partial #1}{\partial #2}}
\newcommand{\Deld}[3]{\dfrac{\partial^{#1} #2}{\partial #3^{#1}}}


\newcommand{\AstIg}{\stackrel{(\ast)}{=}}
\newcommand{\Hop}{\stackrel{L'H\hat{o}pital}{=}}

\newcommand{\red}[1]{{\color{red}#1}} % Para integrales, destacar los cambios.

% Método de integración
\newcommand{\MetInt}[2]{
    \left[\begin{array}{c}
        #1 \\ #2
    \end{array}\right]
}

% Declarar aplicaciones
% 1. Nombre aplicación
% 2. Dominio
% 3. Codominio
% 4. Variable
% 5. Imagen de la variable
\newcommand{\Func}[5]{
    \begin{equation*}
        \begin{array}{rrll}
            \displaystyle #1:& \displaystyle  #2 & \longrightarrow & \displaystyle  #3\\
               & \displaystyle  #4 & \longmapsto & \displaystyle  #5
        \end{array}
    \end{equation*}
}

%------------------------------------------------------------------------


\newcommand{\myhref}[2]{\href{#1}{\color{blue}\uline{#2}}}

\newcommand{\E}{\mathrm{E}}
\newcommand{\V}{\mathrm{Var}}
\newcommand{\R}{\mathbb{R}}
\newcommand{\Rmas}{\R^{+}}
\newcommand{\X}{\mathcal{X}}
\newcommand{\N}{\mathbb{N}}
\newcommand{\I}{\mathrm{Im}}
\newcommand{\Rmenos}{\R^{-}}
\newcommand{\norma}[1]{\lVert #1 \rVert}
\newcommand{\normagenerica}{\lVert \cdot \rVert}
\newcommand{\Rceromas}{\R_0^{+}}
\newcommand{\di}{\text{div}}

\usepackage{pgfplots}
\usepgfplotslibrary{fillbetween}

\begin{document}

    % 1. Foto de fondo
    % 2. Título
    % 3. Encabezado Izquierdo
    % 4. Color de fondo
    % 5. Coord x del titulo
    % 6. Coord y del titulo
    % 7. Fecha

    
    % 1. Foto de fondo
% 2. Título
% 3. Encabezado Izquierdo
% 4. Color de fondo
% 5. Coord x del titulo
% 6. Coord y del titulo
% 7. Fecha
% 8. Autor

\newcommand{\portada}[8]{

    \portadaBase{#1}{#2}{#3}{#4}{#5}{#6}{#7}{#8}
    \portadaBook{#1}{#2}{#3}{#4}{#5}{#6}{#7}{#8}
}

\newcommand{\portadaExamen}[8]{

    \portadaBase{#1}{#2}{#3}{#4}{#5}{#6}{#7}{#8}
    \portadaArticle{#1}{#2}{#3}{#4}{#5}{#6}{#7}{#8}
}




\newcommand{\portadaBase}[8]{

    % Tiene la portada principal y la licencia Creative Commons
    
    % 1. Foto de fondo
    % 2. Título
    % 3. Encabezado Izquierdo
    % 4. Color de fondo
    % 5. Coord x del titulo
    % 6. Coord y del titulo
    % 7. Fecha
    % 8. Autor
    
    
    \thispagestyle{empty}               % Sin encabezado ni pie de página
    \newgeometry{margin=0cm}        % Márgenes nulos para la primera página
    
    
    % Encabezado
    \fancyhead[L]{\helv #3}
    \fancyhead[R]{\helv \nouppercase{\leftmark}}
    
    
    \pagecolor{#4}        % Color de fondo para la portada
    
    \begin{figure}[p]
        \centering
        \transparent{0.3}           % Opacidad del 30% para la imagen
        
        \includegraphics[width=\paperwidth, keepaspectratio]{../../_assets/#1}
    
        \begin{tikzpicture}[remember picture, overlay]
            \node[anchor=north west, text=white, opacity=1, font=\fontsize{60}{90}\selectfont\bfseries\sffamily, align=left] at (#5, #6) {#2};
            
            \node[anchor=south east, text=white, opacity=1, font=\fontsize{12}{18}\selectfont\sffamily, align=right] at (9.7, 3) {\href{https://losdeldgiim.github.io/}{\textbf{Los Del DGIIM}, \texttt{\footnotesize losdeldgiim.github.io}}};
            
            \node[anchor=south east, text=white, opacity=1, font=\fontsize{12}{15}\selectfont\sffamily, align=right] at (9.7, 1.8) {Doble Grado en Ingeniería Informática y Matemáticas\\Universidad de Granada};
        \end{tikzpicture}
    \end{figure}
    
    
    \restoregeometry        % Restaurar márgenes normales para las páginas subsiguientes
    \pagecolor{white}       % Restaurar el color de página
    
    
    \newpage
    \thispagestyle{empty}               % Sin encabezado ni pie de página
    \begin{tikzpicture}[remember picture, overlay]
        \node[anchor=south west, inner sep=3cm] at (current page.south west) {
            \begin{minipage}{0.5\paperwidth}
                \href{https://creativecommons.org/licenses/by-nc-nd/4.0/}{
                    \includegraphics[height=2cm]{../../_assets/Licencia.png}
                }\vspace{1cm}\\
                Esta obra está bajo una
                \href{https://creativecommons.org/licenses/by-nc-nd/4.0/}{
                    Licencia Creative Commons Atribución-NoComercial-SinDerivadas 4.0 Internacional (CC BY-NC-ND 4.0).
                }\\
    
                Eres libre de compartir y redistribuir el contenido de esta obra en cualquier medio o formato, siempre y cuando des el crédito adecuado a los autores originales y no persigas fines comerciales. 
            \end{minipage}
        };
    \end{tikzpicture}
    
    
    
    % 1. Foto de fondo
    % 2. Título
    % 3. Encabezado Izquierdo
    % 4. Color de fondo
    % 5. Coord x del titulo
    % 6. Coord y del titulo
    % 7. Fecha
    % 8. Autor


}


\newcommand{\portadaBook}[8]{

    % 1. Foto de fondo
    % 2. Título
    % 3. Encabezado Izquierdo
    % 4. Color de fondo
    % 5. Coord x del titulo
    % 6. Coord y del titulo
    % 7. Fecha
    % 8. Autor

    % Personaliza el formato del título
    \pretitle{\begin{center}\bfseries\fontsize{42}{56}\selectfont}
    \posttitle{\par\end{center}\vspace{2em}}
    
    % Personaliza el formato del autor
    \preauthor{\begin{center}\Large}
    \postauthor{\par\end{center}\vfill}
    
    % Personaliza el formato de la fecha
    \predate{\begin{center}\huge}
    \postdate{\par\end{center}\vspace{2em}}
    
    \title{#2}
    \author{\href{https://losdeldgiim.github.io/}{Los Del DGIIM, \texttt{\large losdeldgiim.github.io}}
    \\ \vspace{0.5cm}#8}
    \date{Granada, #7}
    \maketitle
    
    \tableofcontents
}




\newcommand{\portadaArticle}[8]{

    % 1. Foto de fondo
    % 2. Título
    % 3. Encabezado Izquierdo
    % 4. Color de fondo
    % 5. Coord x del titulo
    % 6. Coord y del titulo
    % 7. Fecha
    % 8. Autor

    % Personaliza el formato del título
    \pretitle{\begin{center}\bfseries\fontsize{42}{56}\selectfont}
    \posttitle{\par\end{center}\vspace{2em}}
    
    % Personaliza el formato del autor
    \preauthor{\begin{center}\Large}
    \postauthor{\par\end{center}\vspace{3em}}
    
    % Personaliza el formato de la fecha
    \predate{\begin{center}\huge}
    \postdate{\par\end{center}\vspace{5em}}
    
    \title{#2}
    \author{\href{https://losdeldgiim.github.io/}{Los Del DGIIM, \texttt{\large losdeldgiim.github.io}}
    \\ \vspace{0.5cm}#8}
    \date{Granada, #7}
    \thispagestyle{empty}               % Sin encabezado ni pie de página
    \maketitle
    \vfill
}
    \portadaExamen{ffccA4.jpg}{Ecuaciones\\Diferenciales\\ en Mecánica \\ y Biología \\Examen II}{Ecuaciones Diferenciales en Mecánica y Biología. Examen II}{MidnightBlue}{-8}{28}{2026}{José Manuel Sánchez Varbas}

    \begin{description}
        \item[Asignatura] Ecuaciones Diferenciales en Mecánica y Biología.
        \item[Curso Académico] 2023-24.
        \item[Grado] Grado en Matemáticas.
        \item[Descripción] Parcial 1. 
    \end{description}
    \newpage


    % ------------------------------------

    \begin{ejercicio}
        Sea $u_t + t u_x = 0$, con $u : \left[0, T \right[ \times \R \to \R$. Demuestra que $u$ se conserva a lo largo de las trayectorias características.
    \end{ejercicio}

    \begin{ejercicio}
        Resolver 
        \begin{equation*}
            \begin{cases}
                u_t + e^x u_x = 0 \\
                u(0,x) = x
            \end{cases}
        \end{equation*}
    \end{ejercicio}

    \begin{ejercicio}
        Calcular las ondas de choque y rarefacción del problema de Riemann
        $$u_0(x) = \begin{cases}
            -1 \quad& x<0 \\
             1 \quad& x>0
        \end{cases}$$
    \end{ejercicio}

    \newpage
    \setcounter{ejercicio}{0} % Reiniciar contador de ejercicios
    \noindent
    \textbf{Solución.}

    \begin{ejercicio}
        Sea $u_t + t u_x = 0$, con $u : \left[0, T \right[ \times \R \to \R$. Demuestra que $u$ se conserva a lo largo de las trayectorias características. \\

        Obtenemos en primer lugar las trayectorias características. En este caso, $\omega(t,x) = t$. El sistema característico asociado al modelo de transporte es 
        \begin{equation*}
            \begin{cases}
                \dfrac{\partial X}{\partial t}(t;s,x)
                =\omega(t,X(t;s,x)) = t\\[6pt]
                X(s;s,x)=x
            \end{cases}\quad s \in \left[0, T \right[
        \end{equation*}
        %La primera ecuación se resuelve por variables separadas (escribimos $X(t) \equiv X(t;s,x)$):
        %$$\dfrac{dX}{dt} = X' = t \Longrightarrow dX = t dt \Longrightarrow X(t) = \dfrac{t^2}{2} + C \quad C \in \R$$
        %e imponiendo la condición inicial $X(s) = x$:
        %$$\dfrac{s^2}{2} + C = X(s) = x \Longrightarrow C = x - \dfrac{s^2}{2}$$
        %luego
        %$$X(t;s,x) = \dfrac{t^2}{2} + x - \dfrac{s^2}{2} = x + \dfrac{t^2 - s^2}{2}$$
        Para ver que $u$ se conserva a lo largo de las trayectorias características, primeramente hay ver que 
        $$\dfrac{d}{dt}u(t,X(t;s,x)) = 0$$
        Aplicando la regla de la cadena:
        $$
        \dfrac{d}{dt}u(t,X(t;s,x)) =u_t(t,X(t;s,x)) + u_x(t,X(t;s,x))\frac{\partial X}{\partial t}(t;s,x)$$
        Usando que $u$ verifica la ecuación del enunciado, sustituyendo $x$ por $X(t;s,x)$ se tiene que
        $$u_t(t,X(t;s,x))+t\,u_x(t,X(t;s,x))=0$$
        y como 
        $$\dfrac{\partial X}{\partial t}(t;s,x) = \omega(t,X(t;s,x)) = t$$
        por ser una trayectoria característica, nos queda lo que queremos probar:
        $$\dfrac{d}{dt}u(t,X(t;s,x)) = 0$$
        Deducimos entonces que la función $t\mapsto u(t,X(t;s,x))$ es constante, luego
        $$
        u(t,X(t;s,x))=u(s,X(s;s,x)).
        $$

        Y como $X(s;s,x)=x$, nuevamente por ser trayectoria característica, llegamos al resultado
        $$
        \boxed{u(t,X(t;s,x))=u(s,x) \quad \forall t,s \in \left[0, T\right[ \quad \forall x \in \R}
        $$

        Así, el valor de $u$ permanece constante a lo largo de cada trayectoria característica (toma el valor del instante inicial $s$). Por tanto, $u$ se conserva a lo largo de dichas trayectorias, como queríamos demostrar.
    \end{ejercicio}

    \newpage

    \begin{ejercicio}
        Resolver 
        \begin{equation}\label{eq:ec1}
            \begin{cases}
                u_t + e^x u_x = 0 \\
                u(0,x) = x
            \end{cases}
        \end{equation}
        Sabemos por los dos teoremas vistos en clase correspondientes, que tanto la teoría clásica como la teoría débil nos proporcionan existencia y unicidad de soluciones de \eqref{eq:ec1}, y, en caso de verificarse todas las hipótesis, dicha solución única es
        $$u(t,x) = u_0(X(0;t,x))$$
        Sin embargo, en este problema concreto, el teorema para las soluciones clásicas no nos sirve, ya que (H3) no se verifica, puesto que $\omega(t,x) = |e^x|$ no tiene crecimiento sublineal.
        Tampoco podemos aplicar el teorema para las soluciones débiles, puesto que:
        $$\di_x(\omega) = \dfrac{\partial \omega}{\partial x}(t,x) = e^x \neq 0 \quad \forall x \in \R$$
        %Veamos primero las hipótesis para obtener la solución por la vía clásica. En este caso, $\omega(t,x) = e^x$, y $u_0(x) = x$. Primero vemos que se verifican (H1)-(H4):
        %\begin{itemize}
        %    \item[(H1)] Se cumple porque $\omega(t,x)$ es exponencial en la variable $x$, y no depende de $t$. En particular $\omega$ es continua en los dos argumentos.
        %    \item[(H2)] $\omega$ es localmente Lipschitz respecto de $x$, ya que $\omega \in C^1(\left[0,T \right[ \times \R)$, luego $C^1_x(\left[0,T \right[ \times \R)$, y sabemos por un resultado visto en Ecuaciones Diferenciales II que es l.l. resp. de $x$.
        %    Otra forma sería verlo por definición aplicando el Teorema del Valor Medio en un entorno de la forma $I \times J$, con $I = \left]t_0-\delta, t_0 + \delta \right[ \cap \left[0, T\right[$ y $J = \left]x_0-r, x_0 + r \right[$, fijados $(t_0, x_0) \in \left[0, T\right[ \times \R$ y ciertos valores de $\delta, r > 0$.
        %    Así, para cualesquiera $(t,x),(t,y) \in I \times J$, por dicho teorema tendríamos que:
        %    $$|\omega(t,x)-\omega(t,y)| =|e^x-e^y| = e^\xi|x-y| \leq e^{x_0+r}|x-y|$$
        %    para cierto $\xi$ entre $x$ e $y$. Tomando $L = e^{x_0 + r} > 0$ se tiene la condición l.l. resp. de $x$.
        %    \item[(H3)] Como $\omega(t,x) = -tx$, en particular, $$|\omega(t,x)| \leq |t||x| + 0 < T|x| + 0 = A|x| + B$$
        %    para todo $(t,x) \in \left[0, T \right[ \times \R$, por lo que el crecimiento es sublineal.
        %    \item[(H4)] La aplicación \Func{\dfrac{\partial \omega}{\partial x}}{\left[0, T \right[ \times \R}{\R}{(t,x)}{\dfrac{\partial \omega}{\partial x}(t,x) = -t} es polinómica en $t$, luego continua.
        %\end{itemize}
        %Como además es claro que $u_0 \in C^1(\R)$, el teorema mencionado al principio nos dice que el PVI \eqref{eq:ec1} admite una única solución, 
        %$u \in C^1(\left[0, T \right[ \times \R)$, que viene dada por 
        %$$u(t,x) = u_0(X(0;t,x))$$ 
        Aún así, tratamos de buscar una solución a \eqref{eq:ec1} que tenga la misma forma que la que nos proporcionan ambos teoremas, y luego verificar que sea solución. Así, buscamos una solución de la forma que hemos dicho al principio.
        En este caso, $u_0(x) = x$. El sistema característico asociado a \eqref{eq:ec1} es 
        \begin{equation*}
            \begin{cases}
                \dfrac{\partial X}{\partial t}(t;s,x)
                =\omega(t,X(t;s,x)) = e^{X(t;s,x)}\\[6pt]
                X(s;s,x)=x
            \end{cases}\quad s \in \left[0, T \right[
        \end{equation*}
        La primera ecuación se resuelve por variables separadas (escribimos $X(t) \equiv X(t;s,x)$):
        $$\dfrac{dX}{dt} = X' = e^X \Longrightarrow e^{-X}dX = \dfrac{dX}{e^X} = dt \Longrightarrow -e^{-X(t)} = t + C \Longrightarrow$$
        $$e^{-X(t)} = -t - C \Longrightarrow -X(t) = \ln(-t-C) \Longrightarrow X(t) = -\ln(-t-C) \quad C \in \R$$
        e imponiendo la condición inicial $X(s) = x$:
        $$-\ln(-s-C) = X(s) = x \Longrightarrow \ln(-s-C) = -x \Longrightarrow -s - C = e^{-x} \Longrightarrow C = - s - e^{-x}$$
        luego
        $$X(t;s,x) = -\ln(-t-(- s - e^{-x})) = - \ln(e^{-x} + s - t)$$
        Vemos que esta expresión solo está definida cuando el argumento del logaritmo es positivo, es decir,
        $$e^{-x}+s-t>0 \iff t<s+e^{-x}$$
        Esto significa que las trayectorias características no están necesariamente definidas para todo $t\in \left[0,T\right[$, lo cual está relacionado con que no se cumpla (H3).
        Como $X(0;t,x) = - \ln(e^{-x} + t)$, la ``solución'' a la que llegamos de \eqref{eq:ec1} es
        $$\boxed{u(t,x) = u_0(X(0;t,x)) = X(0;t,x) = - \ln(e^{-x} + t)}$$
        Veamos si en efecto es solución. Calculamos:
        $$u_t(t,x)=-\frac{1}{e^{-x}+t}$$
        y
        $$u_x(t,x)=\frac{e^{-x}}{e^{-x}+t}$$
        Luego,
        $$u_t+e^xu_x =-\frac{1}{e^{-x}+t} + \frac{e^xe^{-x}}{e^{-x}+t} =0$$
        Además,
        $$u(0,x)=-\ln(e^{-x})=x$$
        Así, hemos resuelto \eqref{eq:ec1}, como se pedía.
    \end{ejercicio}

    \begin{ejercicio}
        Calcular las ondas de choque y rarefacción del problema de Riemann
        $$u_0(x) = \begin{cases}
            -1 \quad& x<0 \\
             1 \quad& x>0
        \end{cases}$$
        Sabemos por teoría que las soluciones de la ecuación de Burgers (las trayectorias características) son rectas,
        donde la pendiente la determina la condición inicial $u(0,x) = u_0(x)$, o sea, $X(t;0,x) = x + u_0(x)t$. Así pues:
        $$X(t;0,x) = \begin{cases}
            x - t \quad& x < 0 \\
            x + t \quad& x > 0
        \end{cases}$$
        Una solución débil del problema de Riemann (que es el problema de Burgers con un dato inicial con salto) que conocemos es la siguiente:
        $$u_1(t,x) = \begin{cases}
            u^{-} \quad& x < \lm t \\
            u^{+} \quad& x > \lm t
        \end{cases}$$
        donde $u^{-} = -1$ y $u^{+} = 1$ se obtienen del dato inicial. Se ha visto en teoría que el único valor de $\lm$ para 
        que $u$ sea solución (débil) es, aplicando la condición de salto de Rankine-Hugoniot:
        $$\lm = \dfrac{1}{2}(u^{-} + u^{+}) = \dfrac{1}{2}(-1+1) = 0$$
        por lo que ya tenemos una solución del problema de Riemann: 
        $$\boxed{u_1(t,x) = \begin{cases}
            -1 \quad& x < 0 \\
            1 \quad& x > 0
        \end{cases}}$$
        Como además el dato inicial $u_0(x)$ es creciente, sabemos por teoría que tendremos también ondas de rarefacción. Las características ya
        las hemos obtenido:
        $$X(t;0,x) = \begin{cases}
            x - t \quad& x < 0 \\
            x + t \quad& x > 0
        \end{cases}$$
        Obtenemos las rectas que delimitan la región de las ondas de rarefacción haciendo $x \to 0$ en ambos lados, y vemos que:
        $$\lim\limits_{x \to 0^{-}}X(t;0,x) = -t$$
        $$\lim\limits_{x \to 0^{+}}X(t;0,x) = t$$
        Así, las fronteras de la región de rarefacción son $x = -t$ y $x=t$. Las tres regiones que tenemos son $x < -t$, $-t \leq x \leq t$ y $x > t$.
        Según se ha visto en teoría, las ondas de rarefacción pueden construirse como la solución autosemejante o autosimilar $u(t,x) = x/t$,
        y teniendo en cuenta los límites ya obtenidos de la región de rarefacción, una segunda solución (en este caso continua para $t > 0$)
        del problema de este ejercicio es:
        $$\boxed{u_2(t,x) = \begin{cases}
            -1 \quad& x < -t \\
            \dfrac{x}{t} \quad& -t \leq x \leq t \quad t > 0\\
            1 \quad& x > t \\
        \end{cases}}$$
        En $t=0$ tenemos la condición inicial $u(0,x) = u_0(x)$.
    \end{ejercicio}
    
\end{document}
